La baja resolución de una partitura digitalizada puede deteriorar pentagramas, símbolos, dígitos y
texto. La superresolución permite generar una imagen de mayor tamaño, pero los modelos entrenados
con imágenes naturales pueden producir detalles visualmente plausibles que no estaban presentes en
el documento. Este Trabajo de Fin de Grado estudia cuándo estas técnicas aportan valor y qué
precauciones exige su uso profesional en bibliotecas, archivos, conservación, consulta o
preparación editorial.

Se auditó Sheet Music Benchmark (SMB) y se reservó una muestra intacta para la comparación
principal. Esta utilizó 64 páginas de 64 grupos documentales y seis degradaciones sintéticas:
factores de ampliación dos y cuatro, combinados con severidad limpia, moderada y fuerte. Para que la
severidad fuese comparable entre páginas, el desenfoque se normalizó por la separación entre líneas
de pentagrama. Se compararon interpolación bicúbica, la red residual EDSR y el transformador ligero
SwinIR mediante PSNR y SSIM sobre luminancia, intervalos emparejados por grupo, tiempo de inferencia
y una revisión cualitativa preseleccionada de errores de notación. La implementación reproducible
se desarrolló en Python con PyTorch, OpenCV y \texttt{uv}, y la ejecución acelerada se realizó en
Kaggle sobre GPU Tesla T4.

Tras preservar la comparación principal, se ejecutó un estudio secundario de ajuste fino de EDSR.
Se definieron roles internos con 45 grupos de entrenamiento, 13 de validación y 20 de
prueba, sin solapamiento de identificadores ni reutilización de los 64 grupos principales.
Estos grupos pueden ser movimientos de una misma composición: no garantizan independencia
musical y los intervalos son descriptivos, condicionados a esa agrupación. El modelo adaptado superó al
EDSR oficial en PSNR-Y y SSIM-Y para las seis condiciones, con intervalos emparejados que no
cruzaron cero. Este resultado respalda adaptación sobre documentos retenidos de SMB con degradaciones
coincidentes, sin certificar generalización a composiciones nunca vistas ni corrección musical.

Como extensión profesional, el modelo adaptado se aplicó sin reajuste a doce obras de un archivo
musical externo bajo las mismas degradaciones. Los intervalos favorecieron al modelo adaptado en
las seis comparaciones PSNR-Y y en cinco de seis SSIM-Y. Ocho derivados se aceptaron para consulta,
tres con reservas y uno se rechazó; también se implementó un demostrador local reversible. Esta
prueba respalda transferencia entre colecciones bajo LR sintética, no restauración de entradas LR
naturales ni una tasa universal de seguridad.

EDSR y SwinIR superaron a bicúbica en las dos métricas primarias para las seis condiciones. Sin
embargo, SwinIR solo mejoró claramente a EDSR en las condiciones x2 limpia, x2 moderada y x4 limpia,
y fue entre 5,2 y 8,5 veces más lento. La revisión visual detectó símbolos alterados, corrupción de
texto y trazos curvos o texturados incluso cuando aumentaba la fidelidad promedio. Se concluye que
EDSR ofrece el compromiso más defendible entre fidelidad y coste para generar copias de consulta
asistidas, mientras que SwinIR debe reservarse para condiciones donde su ganancia compense el
coste. Ningún método queda validado como restauración semántica, sustituto archivístico o entrada
autónoma de OMR: debe conservarse el original y someterse cada derivado relevante a revisión humana.
